\section*{Hoofdstuk \ref{ch:g12:sums} --- Sommen van toevalsvariabelen en de wet van de grote aantallen}

\begin{solution}{exo:g12:sums:1}
Elke dobbelsteen heeft \omterm{def:g12:randvar:exp}{verwachtingswaarde}
$\frac{1 + 2 + \dots + 6}{6} = \frac{7}{2}$, dus geeft de lineariteit
$\E(S) = \frac72 + \frac72 = 7$. De dobbelstenen zijn \omterm{def:g12:sums:indep}{onafhankelijk}, dus is
$\V(S) = \frac{35}{12} + \frac{35}{12} = \frac{35}{6} \approx 5.83$.
\end{solution}

\begin{solution}{exo:g12:sums:2}
$\E(X) = 50$ en $\V(X) = 100 \times 0.25 = 25$.

Markov ($X \geq 0$): $\P(X \geq 75) \leq \frac{50}{75} = \frac23$.

Bienaymé--Chebyshev: $\P(\abs{X - 50} \geq 25) \leq \frac{25}{25^2} = 0.04$
--- en dat begrenst zelfs de \emph{tweezijdige} \omterm{def:g11:prob:model}{gebeurtenis}, waarvan
$X \geq 75$ de helft is. Chebyshev is hier veel scherper omdat ze de
\omterm{def:g12:randvar:exp}{variantie} uitbuit en niet alleen het \omterm{def:g11:stat:mean}{gemiddelde}. (De echte waarde van
$\P(X \geq 75)$ is kleiner dan $10^{-6}$: beide grenzen zijn grof.)
\end{solution}

\begin{solution}{exo:g12:sums:3}
$F_n$ is het \omterm{prop:g12:sums:mean}{steekproefgemiddelde} van $n$ variabelen van Bernoulli met
parameter $\frac12$ en \omterm{def:g12:randvar:exp}{variantie} $\sigma^2 = \frac14$. De grens
\[
\P\left(\abs{F_n - 0.5} \geq 0.05\right)
\leq \frac{1/4}{n \times 0.05^2} = \frac{100}{n}
\]
is $\leq 0.05$ zodra $n \geq 2000$.
\end{solution}

\begin{solution}{exo:g12:sums:4}
Volgens \cref{prop:g12:sums:mean} is $\E(M_n) = 5\%$ (spreiding verandert het
verwachte rendement niet) en $\sigma(M_n) = \frac{20\%}{\sqrt n}$. De eis
$\frac{20}{\sqrt n} < 4$ geeft $\sqrt n > 5$, dat wil zeggen $n \geq 26$. Dat
is het principe van de \emph{risicospreiding}: kapitaal over onafhankelijke
risico's verdelen deelt het risico (de \omterm{def:g11:stat:variance}{standaardafwijking}) door $\sqrt n$
zonder het verwachte rendement te verlagen.
\end{solution}

\begin{solution}{exo:g12:sums:5}
\emph{1.} Door de $9$ even waarschijnlijke paren te tellen: $S$ neemt de
waarden $2, 3, 4, 5, 6$ aan met de kansen
$\frac19, \frac29, \frac39, \frac29, \frac19$. Dus is
$\E(S) = \frac{2 + 6 + 12 + 10 + 6}{9} = 4$ en
$\E(S^2) = \frac{4 + 18 + 48 + 50 + 36}{9} = \frac{156}{9} = \frac{52}{3}$,
zodat $\V(S) = \frac{52}{3} - 16 = \frac43$.

\emph{2.} Voor één variabele: $\E(X) = 2$,
$\E(X^2) = \frac{1 + 4 + 9}{3} = \frac{14}{3}$ en
$\V(X) = \frac{14}{3} - 4 = \frac23$. Dan is $\E(S) = 2 + 2 = 4$ en, wegens de
\omterm{def:g12:condprob:indep}{onafhankelijkheid}, $\V(S) = \frac23 + \frac23 = \frac43$. Dezelfde waarden.
\end{solution}

\begin{solution}{exo:g12:sums:6}
Met $Y = X$: $\V(X + Y) = \V(2X) = 4\V(X)$, terwijl
$\V(X) + \V(Y) = 2\V(X)$. Beide vallen alleen samen wanneer $\V(X) = 0$, dat
wil zeggen wanneer $X$ constant is --- de optelbaarheid vraagt dus werkelijk
\omterm{def:g12:condprob:indep}{onafhankelijkheid} (hier is $X$ maximaal afhankelijk van zichzelf).
\end{solution}

\begin{solution}{exo:g12:sums:7}
Onder de zuiverheidshypothese is $F_n$ het \omterm{def:g11:stat:mean}{gemiddelde} van $n = 1200$
variabelen van Bernoulli met $p = \frac16$ en
$\sigma^2 = \frac16\cdot\frac56 = \frac{5}{36}$:
\[
\P\left(\abs{F_n - \tfrac16} \geq 0.05\right)
\leq \frac{5/36}{1200 \times 0.0025} = \frac{5}{108} \approx 0.046 .
\]
De waargenomen afwijking is precies $0.05$: een \omterm{def:g11:prob:model}{gebeurtenis} die een zuivere
dobbelsteen met \omterm{def:g10:proba:distribution}{kans} hoogstens $4.6\%$ voortbrengt --- en omdat Chebyshev zeer
behoudend is, ligt de echte \omterm{def:g10:proba:distribution}{kans} veel lager. De zuiverheidshypothese is niet
aannemelijk; de dobbelsteen is zeer waarschijnlijk verzwaard.
\end{solution}

\begin{solution}{exo:g12:sums:8}
\emph{1.} $p(1-p) = \frac14 - \left(p - \frac12\right)^2 \leq \frac14$.

\emph{2.} $F_n$ heeft \omterm{def:g12:randvar:exp}{verwachtingswaarde} $p$ en \omterm{def:g12:randvar:exp}{variantie}
$\frac{p(1-p)}{n} \leq \frac{1}{4n}$; Bienaymé--Chebyshev geeft
\[
\P\bigl(\abs{F_n - p} \geq \varepsilon\bigr)
\leq \frac{p(1-p)}{n\varepsilon^2} \leq \frac{1}{4n\varepsilon^2},
\]
en dat geldt wat de onbekende $p$ ook is.

\emph{3.} Met $\varepsilon = 0.03$ en niveau $0.05$:
\[
\frac{1}{4n(0.03)^2} \leq 0.05
\iff n \geq \frac{1}{4 \times 0.0009 \times 0.05} \approx 5556 .
\]
Ongeveer $5600$ mensen volstaan volgens deze grove grens (het klassieke
antwoord met de normale \omterm{def:g10:numbers:approx}{benadering} ligt dichter bij $1100$, zie
\cref{ch:g12:contdist}).
\end{solution}

\begin{solution}{pb:g12:sums:1}
\textbf{1.} $V(X + Y) = V(X) + V(Y) = \frac{35}{12} + \frac{35}{12} =
\frac{35}{6}$ (\omterm{def:g12:condprob:indep}{onafhankelijkheid}).

\textbf{2.} $\E(\bar X) = 3.5$; $\sigma(\bar X) = \frac{1.71}{\sqrt{100}}
\approx 0.17$: het \omterm{def:g11:stat:mean}{gemiddelde} van honderd dobbelstenen kleeft op een vijfde
van een punt tegen $3.5$ aan.

\textbf{3.} $\P(X \geq 10) \leq \frac{2}{10} = 0.2$.

\textbf{4.} $\P \leq \frac{V(\bar X)}{0.5^2} = \frac{35/1200}{0.25} \approx
0.117$: hoogstens ongeveer $12\,\%$.

\textbf{5.} Met $Y = -X$: $V(X + Y) = V(0) = 0$, terwijl
$V(X) + V(Y) = 2V(X) > 0$: volmaakt tegengesteld gecorreleerde variabelen
heffen elkaar op --- het optellen van \omterm{def:g12:randvar:exp}{varianties} is een voorrecht van de
\omterm{def:g12:condprob:indep}{onafhankelijkheid}.

\textbf{6.} $\E = \frac{18}{37} - \frac{19}{37} = -\frac{1}{37} \approx
-0.027$; $\sigma = \sqrt{1 - \left(\frac{1}{37}\right)^2} \approx 1.00$.

\textbf{7.} $\E(G) = -2.70$ en $\sigma(G) = 10 \times 1.00 = 10$: de
afdrijving bedraagt slechts $0.27\sigma$. De ruis van de avond overstemt het
voordeel --- een grote minderheid van de spelers (ruwweg $40\,\%$, zegt de
klok) stapt met winst naar buiten, en precies dat doet hen terugkomen.

\textbf{8.} De opbrengst van het casino: \omterm{def:g12:randvar:exp}{verwachtingswaarde} $+27\,027$ euro,
\omterm{def:g11:stat:variance}{standaardafwijking} $\approx 1\,000$: de winst ligt $27$ \omterm{def:g11:stat:variance}{standaardafwijkingen}
boven nul. Op $27\sigma$ is ``het casino zou dit jaar kunnen verliezen'' geen
risico, maar een afrondingsfout.

\textbf{9.} $\P(\text{verlies}) = \P(G \geq 27\,027$ voor de spelers$)$
$\leq \frac{10^6 \times 1}{27\,027^2} \approx 0.0014$: zelfs de botste
ongelijkheid uit het boek waarborgt de bank op $99.86\,\%$ --- de waarheid is
astronomisch sterker.

\textbf{10.} Elke ingezette euro heeft, wanneer en hoe ook gekozen, verwacht
rendement $-\frac{1}{37}$ van zichzelf; door de lineariteit is de verwachte
totale winst $-\frac{1}{37} \times (\text{totale inzet})$, negatief voor elke
strategie die iets inzet. Een winnend systeem zou de lineariteit van de
\omterm{def:g12:randvar:exp}{verwachtingswaarde} schenden --- geen enkele slimme volgorde van slechte
inzetten maakt er een goede van. (Verdubbelingssystemen ruilen alleen vele
kleine winsten in voor zeldzame catastrofale verliezen.)

\textbf{11.} De stelling: de \emph{\omterm{def:g10:stats:series}{frequentie}} van rood over $n$ draaien
\omterm{def:g12:seq:limit}{convergeert} (in \omterm{def:g10:proba:distribution}{kans}) naar $\frac{18}{37}$. Ze zegt niets over draai $n + 1$:
het wiel heeft geen geheugen, en na tien keer rood is de \omterm{def:g10:proba:distribution}{kans} op rood nog
altijd $\frac{18}{37}$. Het tegendeel geloven is de \emph{denkfout van de
speler}, en vraag 12 toont wat er werkelijk met reeksen gebeurt.

\textbf{12.} Het overschot van $100$ keer rood wordt niet terugbetaald ---
toekomstige draaien zijn eerlijke kopieën, en het verwachte overschot blijft
$100$. Maar $\frac{100}{10^6} = 0.0001$: een honderdste van een punt. De
vergissing in één zin: \emph{de wet van de grote aantallen verdunt voorbije
toevalligheden in een oceaan van nieuwe proeven; ze stuurt het wiel nooit op
schuldinvordering uit.}

\textbf{13.} $n \geq \frac{1}{4 \times 0.05 \times 0.03^2} \approx 5\,556$
mensen.

\textbf{14.} $\frac{1.96^2 \times 0.25}{0.03^2} \approx 1\,067$: vijf keer
minder. Chebyshev geldt voor \emph{elke} verdeling en betaalt haar
algemeenheid met speling; de constante $1.96$ van de peilers komt van de
werkelijke klokvorm van de schommelingen, die veel harder concentreert.

\textbf{15.} De marge is evenredig met $\frac{1}{\sqrt n}$: ze halveren
verviervoudigt de \omterm{def:g10:proba:sample}{steekproef} --- ongeveer $4\,400$ mensen voor $\pm 1.5$
punten. Precisie wordt tegen kwadratische prijzen gekocht.

\textbf{16.} Het model bestaat uit $n$ onafhankelijke trekkingen met \omterm{def:g10:proba:distribution}{kans} $p$
op succes --- de omvang $N$ van de populatie komt er nooit in voor (een
minuscule fractie van een grote, goed gemengde populatie trekken is precies
wat de \omterm{def:g12:condprob:indep}{onafhankelijkheid} codeert). De \omterm{def:g10:algebra:equation}{vergelijking} van de peilers: om de soep
te proeven volstaat één goed geroerde lepel --- of de pot nu tien of
tienduizend mensen bedient. Het addertje zit in het roeren --- een vertekend
steekproefkader, zoals bij de peiling uit het onderbouwvolume die een heel
land beetnam --- niet in de pot.

\textbf{17.} $315\sqrt n < 30n \iff \sqrt n > 10.5 \iff n > 110$: voorbij
honderd polissen loopt de afdrijving de ruis voorbij, en elke bijkomende polis
vergroot de kloof volgens de wet in $\sqrt n$ --- het bundelen \emph{is} het
bedrijfsmodel.

\textbf{18.} Onafhankelijke gelijke delen:
$\sigma_{\text{pf}} = \frac{20\,\%}{\sqrt{25}} = 4\,\%$: hetzelfde verwachte
rendement en een vijfde van de schommeling. Het maal: risicovermindering
zonder enige kost aan \omterm{def:g12:randvar:exp}{verwachtingswaarde} --- \omterm{def:g12:sums:indep}{onafhankelijk} toeval uitmiddelen.

\textbf{19.} Volmaakt gecorreleerd: de portefeuille is één aandeel in $25$
kostuums: $\sigma = 20\,\%$, geen enkele vermindering. Spreiding huurt de
\emph{\omterm{def:g12:condprob:indep}{onafhankelijkheid}}; wanneer een crisis alles correleert (2008: alle
weddenschappen op vastgoed waren één weddenschap), verdampt de bescherming in
$\sqrt n$ precies wanneer ze nodig is.

\textbf{20.} Sommen drijven af als $n$ en schommelen als $\sqrt n$, zodat de
\omterm{def:g11:stat:mean}{gemiddelden} tot rust komen als $\frac{1}{\sqrt n}$: het casino int de
afdrijving over een miljoen draaien; het peilingsbureau koopt
$\frac{1}{\sqrt{1100}}$ ruis voor een telefoonbudget; de verzekeraar groeit
vanaf $n > 110$ boven zijn eigen schade uit; de belegger \omterm{def:g12:arith:divides}{deelt} zijn risico
door $\sqrt{25}$. De misbruikers: de speler die in schulden gelooft (er is
alleen verdunning), en de financier die in \omterm{def:g12:condprob:indep}{onafhankelijkheid} gelooft (soms is
er maar één weddenschap). De universele vorm van de schommelingen --- de klok
--- is de ster van het continue hoofdstuk, en haar stelling, de centrale
limietstelling, bekroont de kansrekening van de universitaire volumes.
\end{solution}
